% Part: methods
% Chapter: proofs
% Section: reading-proofs

\documentclass[../../../include/open-logic-section]{subfiles}

\begin{document}

\olfileid{mth}{prf}{rea}
\olsection{सिद्धता वाचणे}

पाठ्यपुस्तकांतील आणि लेखांतील सिद्धतांत आपण आत्तापर्यंत उदाहरणांत दिलेला सर्व तपशील क्वचितच असतो. प्रकरणे वेगळी पाहत आहोत, अप्रत्यक्ष सिद्धता देत आहोत किंवा व्याख्या वापरत आहोत, याकडे लेखक अनेकदा स्वतंत्र लक्ष वेधत नाहीत. म्हणून पाठ्यपुस्तकातील सिद्धता वाचताना ती समजण्यासाठी हा तपशील अनेकदा स्वतः भरून काढावा लागतो. सिद्धतेत कराव्या लागणाऱ्या विविध क्रियांचा सराव होण्यासाठीही हे उपयुक्त आहे. एक उदाहरण पाहू.

\begin{prop}[शोषणाचा नियम]
सर्व $A$, $B$ संचांसाठी,
\[
A \cap (A \cup B) = A
\]
\end{prop}

\begin{proof}
$z \in A \cap (A \cup B)$ असेल, तर $z \in A$. म्हणून $A \cap (A \cup B) \subseteq A$. आता $z \in A$ असे समजा. मग $z \in A \cup B$ देखील; म्हणून $z \in A \cap (A \cup B)$ देखील.
\end{proof}

शोषणाच्या नियमाची वरील सिद्धता फार संक्षिप्त आहे. वापरलेल्या व्याख्यांचा उल्लेख नाही; सिद्ध करण्याआधी “हे सिद्ध करायचे आहे” असे म्हटलेले नाही; इत्यादी. ती उलगडू. सिद्ध केलेले विधान हे कोणत्याही $A$ आणि $B$ संचांविषयीचे सामान्य विधान आहे. सिद्धतेत $A$ किंवा $B$ येतात तेव्हा ते स्वेच्छ संचांचे चल आहेत. सिद्धतेत मिळवलेली सामान्य विधाने ही संचांची एकरूपता सिद्ध करण्यासाठी आवश्यक असलेली विधाने आहेत: समीकरणाच्या डाव्या बाजूतील प्रत्येक !!{element} हा उजव्या बाजूतील !!a{element} असतो आणि उलटही तसेच असते.

\begin{quote}
“$z \in A \cap (A \cup B)$ असेल, तर $z \in A$. म्हणून $A \cap (A \cup B) \subseteq A$.”
\end{quote}

हा एकरूपतेच्या सिद्धतेचा पहिला अर्धा भाग आहे. त्यात असे दाखवले आहे: स्वेच्छ~$z$ हा डाव्या बाजूतील !!a{element} असेल, तर तो उजव्या बाजूतील !!a{element} देखील असतो; म्हणजे $A \cap (A \cup B) \subseteq A$. $z \in A \cap (A \cup B)$ असे गृहीत धरा. $z$ हा दोन संचांच्या छेदाचा !!{element} असतो तेव्हाच आणि फक्त तेव्हा जेव्हा तो दोन्ही संचांचा !!{element} असतो. म्हणून $z \in A$ आणि $z \in A \cup B$ देखील, असा निष्कर्ष काढता येतो. विशेषतः $z \in A$ मिळते; हेच दाखवायचे होते. पहिल्या अर्ध्या भागासाठी इतकेच आवश्यक असल्यामुळे उरलेली सिद्धता दुसऱ्या अर्ध्या भागाची असली पाहिजे हे समजते; म्हणजे $A \subseteq A \cap (A \cup B)$ ची सिद्धता.

\begin{quote}
“आता $z \in A$ असे समजा. मग $z \in A \cup B$ देखील; म्हणून $z \in A \cap (A \cup B)$ देखील.”
\end{quote}

सुरुवातीला $z \in A$ असे गृहीत धरतो; कारण कोणत्याही~$z$ साठी $z \in A$ असेल, तर $z \in A \cap (A \cup B)$ असते हे दाखवत आहोत. $z \in A \cap (A \cup B)$ दाखवण्यासाठी “$\cap$” च्या व्याख्येनुसार (i) $z \in A$ आणि (ii) $z \in A \cup B$ दोन्ही दाखवावे लागतात. येथे (i) हे आपले गृहीतकच आहे; म्हणून त्यासाठी आणखी काही सिद्ध करायचे नाही. म्हणूनच सिद्धतेत त्याचा पुन्हा उल्लेख नाही. (ii) साठी हे आठवा: $z$ हा संचांच्या संयोगाचा !!a{element} असतो तेव्हाच आणि फक्त तेव्हा जेव्हा तो त्यांपैकी किमान एका संचाचा !!{element} असतो. येथे $z \in A$ आहे आणि $A \cup B$ हा $A$ आणि $B$ यांचा संयोग आहे; म्हणून ही अट पूर्ण होते. त्यामुळे $z \in A \cup B$. (i) $z \in A$ आणि (ii) $z \in A \cup B$ दोन्ही दाखवले आहेत. म्हणून “$\cap$” च्या व्याख्येनुसार $z \in A \cap (A \cup B)$. सिद्धतेत या व्याख्यांचा उल्लेख नाही; त्या वाचकाला पूर्णपणे परिचित आहेत असे मानलेले आहे. तसे नसल्यास मागे जाऊन त्या काय आहेत हे आठवून घ्यावे लागेल. मग $z \in A$ वरून $z \in A \cup B$ का मिळते आणि $z \in A$ तसेच $z \in A \cup B$ यांवरून $z \in A \cap (A \cup B)$ का मिळते, हेही ओळखावे लागेल.

वरील सिद्धतेचे सर्व तपशील स्पष्ट केलेले आणखी एक रूप असे:
\begin{proof}{}
[संचांसाठीच्या $=$ च्या व्याख्येनुसार, $A \cap (A \cup B) = A$ दाखवण्यासाठी (a) $A \cap (A \cup B) \subseteq A$ आणि (b) $A \subseteq A \cap (A \cup B)$ दाखवावे लागते. (a): $\subseteq$ च्या व्याख्येनुसार $z \in A \cap (A \cup B)$ असेल, तर $z \in A$ असते हे दाखवावे लागते.] $z \in A \cap (A \cup B)$ असेल, तर $z \in A$ [$\cap$ च्या व्याख्येनुसार $z \in A \cap (A \cup B)$ हे $z \in A$ आणि $z \in A \cup B$ असते तेव्हाच आणि फक्त तेव्हाच सत्य असते]. म्हणून $A \cap (A \cup B) \subseteq A$. [(b): $\subseteq$ च्या व्याख्येनुसार $z \in A$ असेल, तर $z \in A \cap (A \cup B)$ असते हे दाखवावे लागते.] आता [(1)] $z \in A$ असे समजा. मग [(2)] $z \in A \cup B$ देखील [(1) वरून $z \in A$ किंवा $z \in B$; आणि $\cup$ च्या व्याख्येनुसार याचा अर्थ $z \in A \cup B$]. म्हणून $z \in A \cap (A \cup B)$ देखील [$\cap$ च्या व्याख्येनुसार $z \in A$, म्हणजे (1), आणि $z \in A \cup B$, म्हणजे (2), दोन्ही आवश्यक आहेत].
\end{proof}

\begin{prob}
$A \cup (A \cap B) = A$ ची पुढील सिद्धता विस्तृत करा. वापरलेल्या सर्व अनुमानाच्या नमुन्यांचा उल्लेख करा; प्रत्येक पायरी गृहीतकांवरून किंवा आधीच्या निष्कर्षांवरून का मिळते आणि कुठे कोणत्या व्याख्यांचा आधार घ्यावा लागतो, हे स्पष्ट करा.
\begin{proof}
$z \in A \cup (A \cap B)$ असेल, तर $z \in A$ किंवा $z \in A \cap B$. $z \in A \cap B$ असेल, तर $z \in A$. कोणताही $z \in A$ हा $\in A \cup (A \cap B)$ देखील असतो.
\end{proof}
\end{prob}

\end{document}
